\jasummary{ライブ壁紙利用を前提とした動的グリッチの表現設計}{}
% TODO: 研究の背景・目的・手法・結果・結論をまとめた和文要旨を記入する。
% テンプレートの例文は転記しない。
スマートフォンの壁紙には、静止画像だけでなく、時間の経過や端末の操作に応じて表示内容が変化する「ライブ壁紙」がある。本研究では、画面の乱れ、色ずれ、ノイズ、部分的な欠損などによって、デジタル機器の表示異常やデータの破損を想起させる「グリッチ表現」に着目した。そして、スマートフォンを傾けることで見た目が変化する、動的かつインタラクティブなグリッチライブ壁紙を実際に制作し、その表現方法と利用者の好みについて検討した。

制作したシステムはAndroid端末を対象とし、OpenGL ESを用いてリアルタイムに描画を行う。元画像をそのまま加工するのではなく、あらかじめ用意した色の組合せや描画設定をもとに、矩形状・ドット状のノイズ、帯状のずれ、RGBずれ、欠損などを画面上に生成し、複数のレイヤとして重ね合わせることでグリッチ表現を構成した。また、端末の傾きを取得するセンサを利用し、傾きに応じて複数のレイヤが異なる量だけ移動する「ゆら」と、端末を動かすことでノイズや色ずれ、ちらつきなどの強度が増加する「ザビ」の2種類の反応方式を実装した。これにより、あらかじめ作成された映像を再生するだけではなく、端末操作に応じてその場で変化するライブ壁紙を実現した。

予備評価では、色彩、ノイズの大きさ、RGBずれ、パララックス、ちらつきなどの設定が異なる8種類のライブ壁紙を6名の参加者に提示し、各表現について「好みではない」から「好み」までの5段階で評価してもらった。当初は、多くの利用者に共通して高く評価されるグリッチ表現を明らかにすることを目的としていたが、評価結果では、参加者によって高く評価する壁紙が異なり、特定の一種類に評価が集中する傾向はみられなかった。

この結果から、グリッチライブ壁紙では、一つの表現を最適なものとして固定するだけではなく、利用者ごとの好みの違いを考慮する必要があると考えた。そこで本研究では、複数の表現から利用者が選択できる構成へと方針を発展させ、動的かつインタラクティブなグリッチ表現に、利用者の選好に応じて表現を選べるパーソナライズ性を取り入れる設計について検討した。
